\section{Тихоновская топология.}
\begin{theoremdefinition}
    Пусть $(X, \tau) \in \cat{Top}$ и $T \neq \emptyset$.
    Тогда тихоновской топологией на $(X, \tau)$ называется слабейшая топология такая, что все проекции $\pi_t\colon X^T \to X, \ x \mapsto x(t)$ являются непрерывными.
    Эта топология имеет предбазу $\sigma_T$ (определённую выше) и, следовательно, совпадает с $\tau_T$.
\end{theoremdefinition}
\begin{proof}
    Поскольку в тихоновской топологии канонические проекции должны быть непрерывными, то есть прообразом всякого открытого должно быть открытое, то $\forall t \in T \ \forall U \in X$ множества вида
    \[
        V(t, U) = \left\{v \in X^T \mid v(t) \in U\right\}.
    \]
    Должны быть открытыми (пусть это будет некоторое семейство $\beta$).
    Заметим, что всякий элемент $\sigma_T$ представим в виде $V(t, U)$, поскольку если есть некоторые $t \in T, \ u \in X^T$ и $O(u(t))$, то
    \[
        W(t, u, O(u(t))) = \left\{v \in X^T \mid v(t) \in O(u(t)) \right\},
    \]
    то $W(t, u, V(u(t)))$ в точности равно $V(t, O)$, а значит $\sigma_T \subset \beta$.
    В тоже время, если $U \neq \emptyset$, то существует $u \in X^T$ такое, что $u(t) \in U$.
    Следовательно $V(t, U) = W(t, u, U(u(t)))$.
    Если $U = \emptyset$, то и $V(t, U) = \emptyset$.
    Значит $\beta = \sigma_T \cup \{\emptyset\}$.
    Пустое множество всегда лежит в топологии, следовательно, топология, порождённая $\sigma_T$ и топология порождённая $\beta$ -- совпадают.
    А значит слабейшая топология, в которой канонические проекции непрерывны, содержит $\tau_T$.

    Поскольку, при этом, в $\tau_T$ проекции очевидно непрерывны, то $\tau_T$ -- искомая топология.
\end{proof}
\begin{notet}
    Доказательство Конста другое.
    Я заменю на него, но чуть позже.
\end{notet}
\begin{example}
    Пусть $T = [0, 1]$ и $\R$ с обычной топологией.
    Пространство $\R^{[0, 1]}$ с тихоновской топологией не является метризуемым, поскольку не удовлетворяет первой аксиоме счётности (а это несложно показать при помощи тех же самых измеримых функций -- секвенциальным замыканием являются они же, а вот замыканием -- всё $\R^{[0, 1]}$).
\end{example}
\begin{example} % дотехать
    Рассмотрим <<пилу>> на отрезке $[0, 1]$, то есть если у нас есть произвольное разбиение $[0, 1]$, то мы во всех точках разбиения принимаем функцию равной нулю, в серединах отрезков разбиения -- единице и соединяем соседние точки прямыми.
    Заметим, что интеграл от всякой <<пилы>> равен $0.5$.
    Следовательно, секвенциальное замыкание множества всех <<пил>> не содержит 0 (в силу теоремы Лебега о мажорируемой сходимости).

    С другой стороны, во всякой окрестности $0$ в топологии Тихонова существует некоторая <<пила>>, поскольку в этой окрестности лежит некоторое множество из базы -- то есть пересечение элементов из предбазы топологии Тихонова.
    Оно имеет вид:
    \[
        W = \left\{x \in \R^{[0, 1]} \mid |x(t_k) - 0| < \epsilon_k\right\}.
    \]
    Тогда $\{t_k\} \cup \{0\} \cup \{1\}$ образуют разбиение отрезка и по нему можно построить пилу.
    А значит замыкание множества <<пил>> содержит $0$.

    Тогда замыкание не совпадает с секвенциальным замыканием и, следовательно, пространство $\R^{[0, 1]}$ не является метризуемым.
\end{example}
\begin{example}
    Построим пример множества такого, что его первое секвенциальное замыкание не совпадает со вторым секвенциальным замыканием.

    Пусть $P$ -- семейство всех неубывающих бесконечно больших последовательностей натуральных чисел.
    Заметим, что $P$ имеет мощность континуума, поскольку оно, в частности, содержит все строго возрастающие бесконечно большие последовательности.
    Каждую такую последовательность $\{n_j\}$ можно представить в виде числа из отрезка $[0, 1]\colon$
    \[
        \alpha = \sum_{j = 1}^{\infty} \frac{1}{2^{n_j}}.
    \]
    Более того, каждое число из отрезка $[0, 1]$ представляется в таком виде.
    А значит $P$ имеет подмножество мощности континуума и, следовательно, само имеет хотя бы континуальную мощность.

    В тоже время, его мощность не более чем континуальна, поскольку если записывать числа в одиннадцатеричной системе, всякую последовательность $\{n_k\}$ можно записать как $0.n_1 A n_2 A \ldots$ -- что и даст нам инъекцию в $[0, 1]$.
    Поэтому $|P| = c$.
    Тогда существует $\phi\colon [0, 1] \to P$ -- некоторая биекция $P$ на отрезок $[0, 1]$.

    Рассмотрим семейство функций $S$, состоящего из
    \[
        S = \left\{ x_{n, m}(t) = \frac{1}{\phi_n(t)} + \frac{n^2}{\phi_m(t)} \mid m, n \in \N \right\},
    \]
    где $\phi_k(t)$ -- функция в $\N$, выдающая на $t$  $k$-й элемент последовательности $\phi(t)$.

    Заметим, что первое секвенциальное замыкание не содержит $0$, а второе -- содержит.

    Очевидно, что второе замыкание содержит $0\colon$
    \[
        \lim_{n \rightarrow \infty} \lim_{m \rightarrow \infty} x_{n, m}(t) = \lim_{n \rightarrow \infty} \frac{1}{\phi_n(t)} = 0.
    \]

    Пусть $0 \in [S]_{seq}$ и существует последовательность $\{n_k, m_k\}$ такая, что
    \[
        x_{n_k, m_k} \rightarrow 0.
    \]
    Без ограничения общности $n_1 < n_2 < \ldots$ (она не может быть ограниченной -- поскольку тогда первый член суммы не стремится к 0; из неограниченной последовательности мы можем выбрать возрастастающую подпоследовательность)
    Также, без ограничения общности, $m_1 < m_2 < \ldots$ (она, аналогично, не может быть ограниченной -- тогда в каждой точке она будет стремится к бесконечности, поскольку $n_k \rightarrow \infty$; далее -- аналогичный аргумент).
    Рассмотрим последовательность вида $x^*(n) = n_1$ при $n \leq m_1$, $x^*(n) = n_2$ при $n \in (m_1, m_2]$ и так далее.
    Заметим, что она лежит в множестве $P$.
    Существует $t^* = \phi^{-1}(x^*)$.
    Но тогда при всяком $m_k\colon$
    \[
        \frac{n_k^2}{\phi_{m_k}(t^*)} = n_k \rightarrow \infty.
    \]
    И поточечная сходимость отсутствует -- значит $0$ действительно не лежит в первом секвенциальном замыкании.
\end{example}